\documentclass[10pt,a4paper]{amsart}
\usepackage[margin=19mm]{geometry}
\usepackage{amsmath,amssymb,mathtools}
\usepackage[scr=rsfs,cal=euler]{mathalfa}
\usepackage{xfrac}
\usepackage[unicode,hidelinks,bookmarks=false]{hyperref}
\newcommand{\ie}{i.e.\ }
\newcommand{\cf}{cf.\ }
\newcommand{\resp}{respectively\ }
\newcommand{\Subsection}{\subsection}
\providecommand{\nobreakdash}{\nobreak\hbox{-}}
\setlength{\emergencystretch}{2em}
\multlinegap0pt
\usepackage{amssymb}
\newtheorem{theorem}{Theorem}
\title{The problem with twp linear branches}
\author{Fritz Schweiger}
\date{Source: arXiv:2602.08493v1 (CC BY 4.0)}
\begin{document}
\maketitle

\noindent\emph{Source and licence.} The problem with twp linear branches. Version arXiv:2602.08493v1 (CC BY 4.0), distributed by arXiv under the Creative Commons Attribution 4.0 International licence, http://creativecommons.org/licenses/by/4.0/, as recorded in the arXiv licence metadata for this exact version and shown on the abstract page https://arxiv.org/abs/2602.08493v1. Full licence notice: \texttt{licenses/arxiv-2602.08493-CC-BY-4.0.txt}.\par\smallskip

\noindent\textit{Editorial scope.} Counted unit: the article's theorem theorem (source0.tex lines 72--77). The article's complete original proof of it is delivered with it (source0.tex lines 79--103, one \texttt{proof} environment of 25 lines). No mathematics is written, repaired or re-derived: the statement and the proof are the article's own lines, in the article's own notation, and no retained line is substituted or re-flowed.  No further source-local context is needed: every cross-reference of every retained line resolves inside the two delivered spans.  Nothing was omitted from any retained span, so the package carries no omission record.  No retained line contains a citation, so the wrapper carries no bibliography and no bibliography entry.  No retained line contains a figure environment, an image-inclusion command or drawing code, so no image is delivered and the package declares no asset dependency.  Editorial formatting only, in addition: an AMS \texttt{amsart} preamble that restates the article's own declarations and macros used by the retained lines, namely the environments \texttt{theorem} on separate counters, together with the standard packages those lines need.  The retained environments print in this document as Theorem 1 in order of appearance, whereas the article prints the same items with the numbering of its own full document.  The complete native source of the article as distributed in the arXiv e-print for this exact version is bundled unchanged as \texttt{sources/194250/source0.tex}.
\begin{theorem} . 
	(1) If the map $T$ is of type $[1,1,1]$, then  $S$ has a natural dual and 
	$$ h(x) = \frac{1}{(2 - \beta + 3 \beta x) (2 + 3 \beta x )}.$$
	(1) If the map $T$ is of type $[1,-1,1]$, then  $S$ has a natural dual and 
$$ h(x) = \frac{1}{(4+ \beta + 3 \beta x) (4 + 3 \beta x )}.$$
\end{theorem}

	\begin{proof}
	We list the relevant (inverse) branches of $T$ and $S$.
	$$ \varepsilon_\alpha =1: V_\alpha x = \frac{-1+3x}{3}  , , \varepsilon_\alpha =-1: V_\alpha x = \frac{2 - 3x}{3} $$          
		$$ \varepsilon_\beta =1: V_\beta x = \frac{1 + (1+2 \beta)x}{3 + 3 \beta x}, \, \varepsilon_\beta = -1 : V_\beta x = \frac{2 + (-1 + \beta )x}{3 + 3\beta x}$$
		$$ \varepsilon_\gamma =1: V_\gamma x = \frac{1+3x}{3}  , , \varepsilon_\gamma =-1: V_\gamma x = \frac{4 - 3x}{3}           $$
	To determine the map $S$ one calculates  $V_{\alpha \beta}x = V_\alpha (V_\beta x)$
	and $V_{\gamma \beta}x = V_\gamma (V_\beta x)$.\\	
	We obtain
	$$ \varepsilon_{\alpha\beta} = 1: V_{\alpha\beta} x = \frac{(1+\beta)x}{3 + 3 \beta x}, \,\varepsilon_{\alpha\beta} = -1: V_{\alpha\beta} x = \frac{(1 - x}{3 + 3 \beta x}$$
		$$ \varepsilon_{\gamma\beta} = 1: V_{\gamma\beta} x = \frac{2+(1+3\beta)x}{3 + 3 \beta x}, \,\varepsilon_{\gamma\beta} = -1: V_{\gamma\beta} x = \frac{3 +(-1+2\beta)  x}{3 + 3 \beta x}$$
	 We look for a map $Mx = \frac{B + D x}{A+ Bx}$ such that $M (V_i x) = V_i^{*} (Mx)$, $ i \in \{\alpha \beta, \beta, \gamma \beta\}$.
	Then we obtain three homogeneous equations for three indeterminants $A, B, D$. A non-trivial  soltion requires that the Determinant vanishes,  $DET = 0$.\\ 
	(1)  Type $[1,1,1]$\\
	We find the homogeneous equations
	$$ A(3\beta) + B(-2 + \beta) = 0$$
	$$ A 3\beta + B(-2 + 2\beta) - D =0$$
	$$A3\beta + B (-2 + 3\beta) -2D = 0.$$
	Then $DET = 0$ and we find $A= 2- \beta$, $B= 3\beta$, and $D=3\beta^2$.\\
	As an example we take $\beta = 2$. Then we have $$V_{\alpha \beta}x  = \frac{x}{1+2x}, V_\beta x = \frac{1+5x}{3+6x}, V_{\gamma \beta}x = \frac{2+7x}{3+6x} $$ and $$h(x) = \frac{1}{x(1+3x)}.$$
	
	(2) Type $[1,-1,1]$\\
	Again we find $DET = 0$ and the values $A= 4 + \beta$, $B= 3 \beta$, and $D= 3 \beta^2$.	
		As an example we take $\beta = 1$. Then we have $$V_{\alpha \beta}x  = \frac{1-x}{3+x}, V_\beta x = \frac{2}{3+3x}, V_{\gamma \beta}x = \frac{3+x}{3+3x} $$ and $$h(x) = \frac{1}{(4+3x)(5+3x)}.$$
		
	\end{proof}

\end{document}
